\section{Conclusion and future work}\label{sec:Conclusion and future work}

\paragraph{Conclusion.}
This thesis has developed a rigorous probabilistic framework for finite-context autoregressive transformers acting on sequences of functions. The proposed architecture combines a bounded function-space transformer with coefficient-valued output heads in a fixed approximation space, producing a context-dependent Gaussian law with non-diagonal covariance and a positive covariance floor. Taking the complete functional context as the state turns recursive shift-and-append generation into a time-homogeneous general-state-space Markov chain. The conditional Gaussian laws are shown to be well defined, to have common support, and to depend measurably on the context. Their common support and covariance bounds yield a global \(n\)-step Doeblin minorisation, which gives a unique invariant distribution and uniform geometric convergence in total variation from every initial distribution; the Foster--Lyapunov drift estimate separately provides finite weighted second-moment control. For a fixed parameter value, we also establish two dependent-data pretraining guarantees: a Marton-coupling bound for bounded Wasserstein transition losses and a transportation-based bound for Lipschitz losses under metric dependence, both with explicit sample-complexity statements under zero empirical risk. The analysis therefore includes probability-kernel well-posedness, approximation restrictions, long-time inference behaviour, and fixed-parameter pretraining generalisation.

The in-context analysis extends these concentration tools to fixed pretrained models evaluated on functional task trajectories. A model-comparison lemma for squared Wasserstein risk gives simultaneous comparator bounds, explicit mixing-time and metric dependence factors, and sample complexities above a residual task error. Heat-equation trajectories and a recursion with infinite-rank functional noise verify the assumptions in concrete settings. Projection error and the positive covariance floor quantify restrictions on attainable accuracy, while the distinction between evaluated windows and visible context length separates risk concentration from improvement with longer prompts.

\paragraph{Future work.}
Several directions follow from the results of this thesis. First, the approximation--stability trade-offs associated with the approximation dimension, covariance floor, context length, and conditional output law should be studied theoretically and empirically. Adaptive or increasing approximation spaces could reduce projection error, while Gaussian mixtures or conditional function-space diffusion mechanisms may provide more expressive transition laws. Second, the global boundedness, common-support, and covariance-floor conditions used to obtain the Doeblin minorisation could be relaxed. Models admitting only local small sets may instead require Harris-type drift conditions, while sharper arguments would be needed to obtain quantitatively useful convergence rates. Third, empirical experiments should assess recursive generation, sensitivity to the architectural parameters appearing in the bounds, and the extent to which the predicted long-run behaviour is observed. On the statistical side, an important extension is to move beyond fixed-parameter concentration and control parameters selected from the same dependent data, for example through uniform or algorithm-dependent arguments adapted to functional observations. Finally, combining temporal memory with functional demonstrations would connect the present autoregressive setting with in-context operator learning, allowing an inferred operator to be used for uncertainty-aware recursive generation.
